\section{AMC 12 12B 2004}
        \begin{enumerate}
        	\item (\textit{AMC 12 12B 2004}) At each basketball practice last week, Jenny made twice as many free throws as she made at the previous practice. At her fifth practice she made 48 free throws. How many free throws did she make at the first practice? 


 $(\mathrm {A}) 3\qquad (\mathrm {B}) 6 \qquad (\mathrm {C}) 9 \qquad (\mathrm {D}) 12 \qquad (\mathrm {E}) 15$

 \textbf{Answer: }A



\textbf{Solution}

Each day Jenny makes half as many free throws as she does at the next practice. Hence on the fourth day she made $\frac{1}{2} \cdot 48 = 24$ free throws, on the third $12$, on the second $6$, and on the first $3 \Rightarrow \mathrm{(A)}$.

 Because there are five days, or four transformations between days (day 1 $\rightarrow$ day 2 $\rightarrow$ day 3 $\rightarrow$ day 4 $\rightarrow$ day 5), she makes $48 \cdot \frac{1}{2^4} = \boxed{\mathrm{(A)}\ 3}$



	\item (\textit{AMC 12 12B 2004}) In the expression $c\cdot a^b-d$, the values of $a$, $b$, $c$, and $d$ are 0, 1, 2, and 3, although not necessarily in that order. What is the maximum possible value of the result?


 $(\mathrm {A}) 5\qquad (\mathrm {B}) 6 \qquad (\mathrm {C}) 8 \qquad (\mathrm {D}) 9 \qquad (\mathrm {E}) 10$

 \textbf{Answer: }D



\textbf{Solution 1}

If $a=0$ or $c=0$, the expression evaluates to $-d<0$. 
If $b=0$, the expression evaluates to $c-d\leq 2$. 
Case $d=0$ remains.In that case, we want to maximize $c\cdot a^b$ where $\{a,b,c\}=\{1,2,3\}$. Trying out the six possibilities we get that the greatest is $(a,b,c)=(3,2,1)$, where $c\cdot a^b=1\cdot 3^2=\boxed{\mathrm{(D)}\ 9}$.

 



\textbf{Solution 2}

We first minimize $d$ as $0$, since is it subtracted and is independent from the rest. Then we are left with $a$, $b$, $c$ is from the set ${1,2,3}$. Then after some logical tries we can get $1\cdot3^2-0=\boxed{\mathrm{(D)}\ 9}$ is the largest.

 


	\item (\textit{AMC 12 12B 2004}) If $x$ and $y$ are positive integers for which $2^x3^y=1296$, what is the value of $x+y$?


 $(\mathrm {A}) 8\qquad (\mathrm {B}) 9 \qquad (\mathrm {C}) 10 \qquad (\mathrm {D}) 11 \qquad (\mathrm {E}) 12$

 \textbf{Answer: }A



\textbf{Solution}

$1296 = 2^4 3^4$ and $4+4=\boxed{8} \Longrightarrow \mathrm{(A)}$.

 


	\item (\textit{AMC 12 12B 2004}) An integer $x$, with $10\leq x\leq 99$, is to be chosen. If all choices are equally likely, what is the probability that at least one digit of $x$ is a 7?


 $(\mathrm {A}) \dfrac{1}{9} \qquad (\mathrm {B}) \dfrac{1}{5} \qquad (\mathrm {C}) \dfrac{19}{90} \qquad (\mathrm {D}) \dfrac{2}{9} \qquad (\mathrm {E}) \dfrac{1}{3}$

 \textbf{Answer: }B



\textbf{Solution}

The digit $7$ can be either the tens digit ($70, 71, \dots, 79$: $10$ possibilities), or the ones digit ($17, 27, \dots, 97$: $9$ possibilities), but we counted the number $77$ twice. This means that out of the $90$ two-digit numbers, $10+9-1=18$ have at least one digit equal to $7$. Therefore the probability is $\dfrac{18}{90} = \boxed{\dfrac{1}{5}} \Longrightarrow \mathrm{(B)}$.

 



\textbf{Solution 2}

By complementary counting, we count the numbers that do not contain a $7$, then subtract from the total. There is a $\frac{8}{9}\cdot\frac{9}{10}$ probability of choosing a number that does NOT contain a $7$. Subtract this from $1$ and simplify yields $1 - \frac{8}{9}\cdot\frac{9}{10} = \frac{90}{90} - \frac{72}{90} = \frac{18}{90} = \frac{1}{5}$.

 


	\item (\textit{AMC 12 12B 2004}) On a trip from the United States to Canada, Isabella took $d$ U.S. dollars. At the border she exchanged them all, receiving 10 Canadian dollars for every 7 U.S. dollars. After spending 60 Canadian dollars, she had $d$ Canadian dollars left. What is the sum of the digits of $d$? 


 $(\mathrm {A}) 5\qquad (\mathrm {B}) 6 \qquad (\mathrm {C}) 7 \qquad (\mathrm {D}) 8 \qquad (\mathrm {E}) 9$

 \textbf{Answer: }A



\textbf{Solution 1}

Isabella had $60+d$ Canadian dollars. Setting up an equation we get $d=\frac{7}{10}\cdot(60+d)$, which solves to $d=140$, and the sum of digits of $d$ is $\boxed{\mathrm{(A)}\ 5}$.

 



\textbf{Solution 2}

Each time Isabella exchanges $7$ U.S. dollars, she gets $7$ Canadian dollars and $3$ Canadian dollars extra. Isabella received a total of $60$ Canadian dollars extra, therefore she exchanged $7$ U.S. dollars $\frac{60}{3}=20$ times. Thus $d=7\cdot20=140$, and the sum of the digits is $\boxed{\mathrm{(A)}\ 5}$.

 



\textbf{Solution 3}

Isabella started with $d$ dollars and ended with $d$ dollars. Therefore,

 $\frac{d}{7}\cdot10-60=d$.

 Solving for $d$, we obtain $d = 140$. The sum of the digits of $d$ is $\boxed{\mathrm{(A)}\ 5}$.

 


	\item (\textit{AMC 12 12B 2004}) Minneapolis-St. Paul International Airport is 8 miles southwest of downtown St. Paul and 10 miles southeast of downtown Minneapolis. Which of the following is closest to the number of miles between downtown St. Paul and downtown Minneapolis? 


 $(\mathrm {A}) 13\qquad (\mathrm {B}) 14 \qquad (\mathrm {C}) 15 \qquad (\mathrm {D}) 16 \qquad (\mathrm {E}) 17$

 \textbf{Answer: }A



\textbf{Solution}

The directions "southwest" and "southeast" are orthogonal. Thus the described situation is a right triangle with legs $8$ miles and $10$ miles long. The hypotenuse length is $\sqrt{8^2 + 10^2}\approx12.8$, and thus the answer is $\boxed{\mathrm{(A)}\ 13}$.

 Without a calculator one can note that $8^2+10^2=164<169=13^2\Rightarrow\boxed{\mathrm{(A)}\ 13}$.

 


	\item (\textit{AMC 12 12B 2004}) A square has sides of length 10, and a circle centered at one of its vertices has radius 10. What is the area of the union of the regions enclosed by the square and the circle?


 $(\mathrm {A}) 200+25\pi \quad (\mathrm {B}) 100+75\pi \quad (\mathrm {C}) 75+100\pi \quad (\mathrm {D}) 100+100\pi \quad (\mathrm {E}) 100+125\pi$

 \textbf{Answer: }B



\textbf{Solution}

The area of the circle is $S_{\bigcirc}=100\pi$; the area of the square is $S_{\square}=100$.

 Exactly $\frac{1}{4}$ of the circle lies inside the square. Thus the total area is $\dfrac34 S_{\bigcirc}+S_{\square}=\boxed{\mathrm{(B)\ }100+75\pi}$.

 
                    \begin{figure}[H]
                        \centering
                        \includegraphics[width=0.5\textwidth]{Solution/AMC_12_12B_2004_7_Solution_Figure_1.png}
                    \end{figure}
                    



	\item (\textit{AMC 12 12B 2004}) A grocer makes a display of cans in which the top row has one can and each lower row has two more cans than the row above it. If the display contains 100 cans, how many rows does it contain?


 $(\mathrm {A}) 5 \qquad (\mathrm {B}) 8 \qquad (\mathrm {C}) 9 \qquad (\mathrm {D}) 10 \qquad (\mathrm {E}) 11$

 \textbf{Answer: }D



\textbf{Solution}

The sum of the first $n$ odd numbers is $n^2$. As in our case $n^2=100$, we have $n=\boxed{\mathrm{(D)\ }10}$.

 


	\item (\textit{AMC 12 12B 2004}) The point $(-3,2)$ is rotated $90^\circ$ clockwise around the origin to point $B$. Point $B$ is then reflected over the line $x=y$ to point $C$. What are the coordinates of $C$?


 $\mathrm{(A)}\ (-3,-2) \qquad \mathrm{(B)}\ (-2,-3) \qquad \mathrm{(C)}\ (2,-3) \qquad \mathrm{(D)}\ (2,3) \qquad \mathrm{(E)}\ (3,2)$

 \textbf{Answer: }E



\textbf{Solution}

The entire situation is in the picture below. The correct answer is $\boxed{\mathrm{(E)}\ (3,2)}$.

 
                    \begin{figure}[H]
                        \centering
                        \includegraphics[width=0.5\textwidth]{Solution/AMC_12_12B_2004_9_Solution_Figure_1.png}
                    \end{figure}
                    



	\item (\textit{AMC 12 12B 2004}) An annulus is the region between two concentric circles. The concentric circles in the figure have radii $b$ and $c$, with $b>c$. Let $OX$ be a radius of the larger circle, let $XZ$ be tangent to the smaller circle at $Z$, and let $OY$ be the radius of the larger circle that contains $Z$. Let $a=XZ$, $d=YZ$, and $e=XY$. What is the area of the annulus? 


 
                    \begin{figure}[H]
                        \centering
                        \includegraphics[width=0.5\textwidth]{Problem/AMC_12_12B_2004_10_Problem_Figure_1.png}
                    \end{figure}
                    
$\mathrm{(A) \ } \pi a^2 \qquad \mathrm{(B) \ } \pi b^2 \qquad \mathrm{(C) \ } \pi c^2 \qquad \mathrm{(D) \ } \pi d^2 \qquad \mathrm{(E) \ } \pi e^2$

 \textbf{Answer: }A



\textbf{Solution}

The area of the large circle is $\pi b^2$, the area of the small one is $\pi c^2$, hence the shaded area is $\pi(b^2-c^2)$.

 From the \href{https://artofproblemsolving.com/wiki/index.php?title=Pythagorean\_Theorem}{Pythagorean Theorem} for the right triangle $OXZ$ we have $a^2 + c^2 = b^2$, hence $b^2-c^2=a^2$ and thus the shaded area is $\boxed{\mathrm{(A)\ }\pi a^2}$.

 



\textbf{Solution 2}

Set $c=0,$ then the shaded area is just the area of a circle with radius $a,$ which is $\boxed{\mathrm{(A)\ }\pi a^2}$.

 


	\item (\textit{AMC 12 12B 2004}) All the students in an algebra class took a $100$-point test. Five students scored $100$, each student scored at least $60$, and the mean score was $76$. What is the smallest possible number of students in the class?


 $\mathrm{(A)}\ 10 \qquad \mathrm{(B)}\ 11 \qquad \mathrm{(C)}\ 12 \qquad \mathrm{(D)}\ 13 \qquad \mathrm{(E)}\ 14$

 \textbf{Answer: }D



\textbf{Solution 1}

Let the number of students be $n\geq 5$. Then the sum of their scores is at least $5\cdot 100 + (n-5)\cdot 60$. At the same time, we need to achieve the mean $76$, which is equivalent to achieving the sum $76n$.

 Hence we get a necessary condition on $n$: we must have $5\cdot 100 + (n-5)\cdot 60 \leq 76n$.This can be simplified to $200 \leq 16n$. The smallest integer $n$ for which this is true is $n=13$.

 To finish our solution, we now need to find one way how $13$ students could have scored on the test. We have $13\cdot 76 = 988$ points to divide among them. The five $100$s make $500$, hence we must divide the remaining $488$ points among the other $8$ students. This can be done e.g. by giving $61$ points to each of them. 

 Hence the smallest possible number of students is $\boxed{\mathrm{(D)}\ 13}$.

 



\textbf{Solution 2}

Let's assume that each student scores 60 points on the test. This gives$(500 + 60(n))/(n+5) = 76$

 Solving the equation for $n$ we see that $n = 7.5$.

 Because of the extra five people we know that the total amount of people is $12.5$, however, to adjust for the complexity of not every score being a $60$ we round up to $13$. Giving $\boxed{\mathrm{(D)}\ 13}$.

 


	\item (\textit{AMC 12 12B 2004}) In the sequence $2001$, $2002$, $2003$, $\ldots$ , each term after the third is found by subtracting the previous term from the sum of the two terms that precede that term. For example, the fourth term is $2001 + 2002 - 2003 = 2000$. What is the
$2004^\textrm{th}$ term in this sequence?


 $\mathrm{(A) \ } -2004 \qquad \mathrm{(B) \ } -2 \qquad \mathrm{(C) \ } 0 \qquad \mathrm{(D) \ } 4003 \qquad \mathrm{(E) \ } 6007$

 \textbf{Answer: }C



\textbf{Solution 1}

We already know that $a_1=2001$, $a_2=2002$, $a_3=2003$, and $a_4=2000$. Let's compute the next few terms to get the idea how the sequence behaves. We get $a_5 = 2002+2003-2000 = 2005$, $a_6=2003+2000-2005=1998$, $a_7=2000+2005-1998=2007$, and so on.

 We can now discover the following pattern: $a_{2k+1} = 2001+2k$ and $a_{2k}=2004-2k$. This is easily proved by induction. It follows that $a_{2004}=a_{2\cdot 1002} = 2004 - 2\cdot 1002 = \boxed{0}$.

 



\textbf{Solution 1 but in a more not smart way}

Subtract 2000 from each of the terms so the sequence turns into 1, 2, 3, 0, 5, -2, 7, -4, 9, -6... (1+2-3=0, 2+3-0=5, 3+0-5=-2, 0+5-(-2)=7, 5+(-2)-7=-4, etc.). Quickly notice that after taking away 1 and 2, the 2nd term is 0, the 4th term is -2, the 6th term is -4, the 8th term is -6, etc. Thus, in the sequence without 1 and 2, the (2n)th term has a value of -2(n-1). Therefore, to find the 2004th term in the original sequence, we take away the first 2 values to form the new sequence so the value of the 2004th term in the original sequence is the 2002nd term in the new sequence. 2002 is 2 times 1001. Thus, the value is -2(1001-1)= -2000. BUT, REMEMBER THAT AT THE START WE TOOK AWAY 2000 FROM EACH TERM. ADD -2000 TO 2000 TO GET $\boxed{0}$. LOLOL

 



\textbf{Solution 2}

Note that the recurrence $a_n+a_{n+1}-a_{n+2}~=~a_{n+3}$ can be rewritten as $a_n+a_{n+1} ~=~ a_{n+2}+a_{n+3}$. 

 Hence we get that $a_1+a_2 ~=~ a_3+a_4 ~=~ a_5+a_6 ~= \cdots$ and also$a_2+a_3 ~=~ a_4+a_5 ~=~ a_6+a_7 ~= \cdots$

 From the values given in the problem statement we see that $a_3=a_1+2$. 

 From $a_1+a_2 = a_3+a_4$ we get that $a_4=a_2-2$.

 From $a_2+a_3 = a_4+a_5$ we get that $a_5=a_3+2$.

 Following this pattern, we get $a_{2004} = a_{2002} - 2 = a_{2000} - 4 = \cdots = a_2 - 2002 = \boxed{0}$.

 



\textbf{Solution 3}

Our recurrence is $a_n+a_{n+1}-a_{n+2}~=~a_{n+3}$, so we get $r^3+r^2-r-1 = 1$, so $(r-1)(r+1)^2 = 1$, so our formula for the recurrence is $a_n = A + (B + Cn)(-1)^n$.

 Substituting our starting values gives us $a_n = 2002 + (2 - n)(-1)^n$.

 So, $a_{2004} = 2002 - 2002 = 0.$



	\item (\textit{AMC 12 12B 2004}) If $f(x) = ax+b$ and $f^{-1}(x) = bx+a$ with $a$ and $b$ real, what is the value of $a+b$? 


 $\mathrm{(A)}\ -2 \qquad\mathrm{(B)}\ -1 \qquad\mathrm{(C)}\ 0 \qquad\mathrm{(D)}\ 1 \qquad\mathrm{(E)}\ 2$

 \textbf{Answer: }A



\textbf{Solution (Alcumus)}

Since $f(f^{-1}(x))=x$, it follows that $a(bx+a)+b=x$, which implies $abx + a^2 +b = x$. This equation holds for all values of $x$ only if $ab=1$ and $a^2+b=0$. 

 Then $b = -a^2$. Substituting into the equation $ab = 1$, we get $-a^3 = 1$. Then $a = -1$, so $b = -1$, and
 \[f(x)=-x-1.\]Likewise
 \[f^{-1}(x)=-x-1.\]These are inverses to one another since
 \[f(f^{-1}(x))=-(-x-1)-1=x+1-1=x.\]
 \[f^{-1}(f(x))=-(-x-1)-1=x+1-1=x.\]Therefore $a+b=\boxed{\mathrm{(A)}\ -2}$.

 note: you could just find that $a=-1$ and $b=-1$ and add, getting $-2$.

 


	\item (\textit{AMC 12 12B 2004}) In $\triangle ABC$, $AB=13$, $AC=5$, and $BC=12$. Points $M$ and $N$ lie on $AC$ and $BC$, respectively, with $CM=CN=4$. Points $J$ and $K$ are on $AB$ so that $MJ$ and $NK$ are perpendicular to $AB$. What is the area of pentagon $CMJKN$? 


 
                    \begin{figure}[H]
                        \centering
                        \includegraphics[width=0.5\textwidth]{Problem/AMC_12_12B_2004_14_Problem_Figure_1.png}
                    \end{figure}
                    
$\mathrm{(A)}\ 15 \qquad \mathrm{(B)}\ \frac{81}{5} \qquad \mathrm{(C)}\ \frac{205}{12} \qquad \mathrm{(D)}\ \frac{240}{13} \qquad \mathrm{(E)}\ 20$

 \textbf{Answer: }D



\textbf{Solution 1}

The triangle $ABC$ is clearly a right triangle, its area is $\frac{5\cdot 12}2 = 30$. If we knew the areas of triangles $AMJ$ and $BNK$, we could subtract them to get the area of the pentagon.

 Draw the height $CL$ from $C$ onto $AB$. As $AB=13$ and the area is $30$, we get $CL=\frac{60}{13}$. The situation is shown in the picture below:

 
                    \begin{figure}[H]
                        \centering
                        \includegraphics[width=0.5\textwidth]{Solution/AMC_12_12B_2004_14_Solution_Figure_1.png}
                    \end{figure}
                    

 Now note that the triangles $ABC$, $AMJ$, $ACL$, $CBL$ and $NBK$ all have the same angles and therefore they are similar. We already know some of their sides, and we will use this information to compute their areas. Note that if two polygons are similar with ratio $k$, their areas have ratio $k^2$. We will use this fact repeatedly.Below we will use $[XYZ]$ to denote the area of the triangle $XYZ$.

 We have $\frac{CL}{BC} = \frac{60/13}{12} = \frac 5{13}$, hence $[ACL] = \frac{ 25[ABC] }{169} = \frac{750}{169}$.

 Also, $\frac{CL}{AC} = \frac{60/13}5 = \frac{12}{13}$, hence $[CBL] = \frac{ 144[ABC] }{169} = \frac{4320}{169}$.

 Now for the smaller triangles:

 We know that $\frac{AM}{AC}=\frac 15$, hence $[AMJ] = \frac{[ACL]}{25} = \frac{30}{169}$. 

 Similarly, $\frac{BN}{BC}=\frac 8{12} = \frac 23$, hence $[NBK] = \frac{4[CBL]}9 = \frac{1920}{169}$.

 Finally, the area of the pentagon is $30 - \frac{30}{169} - \frac{1920}{169} = \boxed{\frac{240}{13}}$.

 



\textbf{Solution 1a}

Observe that all of the triangles in the problem are right triangles and similar with a ratio of 5-12-13. The largest triangle $\Delta ABC$ has area $\dfrac{5\cdot12}2=30$. $\Delta AMJ$ is linearly scaled down from $\Delta ABC$ by a factor of $\dfrac{5-4}{13}=\dfrac1{13}$ (as we can see from comparing the two hypotenuses), while $\Delta NBK$ is scaled by a factor of $\dfrac{12-4}{13}=\dfrac8{13}$. The area we desire is the combined area of $\Delta AMJ$ and $\Delta NBK$ subtracted from the area of $\Delta ABC$, which is 

 
 \[30-30\left(\dfrac1{13}\right)^2-30\left(\dfrac8{13}\right)^2=30\left(1-\left(\dfrac1{13}\right)^2-\left(\dfrac8{13}\right)^2\right)=30\left(1-\dfrac{65}{169}\right)=30\left(\dfrac8{13}\right)=\boxed{\textbf{(D) }\dfrac{240}{13}}.\]




\textbf{Solution 2}

Split the pentagon along a different diagonal as follows:

 
                    \begin{figure}[H]
                        \centering
                        \includegraphics[width=0.5\textwidth]{Solution/AMC_12_12B_2004_14_Solution_Figure_2.png}
                    \end{figure}
                    The area of the pentagon is then the sum of the areas of the resulting right triangle and trapezoid. As before, triangles $ABC$, $AMJ$, and $NBK$ are all similar.

 Since $BN=12-4=8$, $NK=\frac{5}{13}(8)=\frac{40}{13}$ and $BK=\frac{12}{13}(8)=\frac{96}{13}$. Since $AM=5-4=1$, $JM=\frac{12}{13}$ and $AJ=\frac{5}{13}$.

 The trapezoid's height is therefore $13-\frac{5}{13}-\frac{96}{13}=\frac{68}{13}$, and its area is $\frac{1}{2}\left(\frac{68}{13}\right)\left(\frac{12}{13}+\frac{40}{13}\right)=\frac{34}{13}(4)=\frac{136}{13}$.

 Triangle $MCN$ has area $\frac{1}{2}(4)(4)=8$, and the total area is $\frac{104+136}{13}=\boxed{\frac{240}{13}}$.

 



\textbf{Solution 3}

Right triangle $ABC$ has area $1/2 \cdot AC \cdot BC = 1/2 \cdot 5 \cdot 12 = 30$.

 We are given that $AC = 5$ and $CM = 4$, so $AM = AC - CM = 1$. Right triangles $AMJ$ and $ABC$ are similar, with hypotenuses 1 and 13, respectively, so triangle $AMJ$ has area
 \[\left( \frac{1}{13} \right)^2 \cdot 30.\]Also, we are given that $BC = 12$ and $CN = 4$, so $BN = BC - CN = 8$. Right triangles $NBK$ and $ABC$ are similar, with hypotenuses 8 and 13, respectively, so triangle $NBK$ has area
 \[\left( \frac{8}{13} \right)^2 \cdot 30.\]Therefore, the area of pentagon $CMJKN$ is
 \[30 - \left( \frac{1}{13} \right)^2 \cdot 30 - \left( \frac{8}{13} \right)^2 \cdot 30 = \boxed{\frac{240}{13}}.\]




\textbf{Solution 4}

Because triangle ABC, triangle NBK, and triangle AMJ are similar right triangles whose hypotenuses are in the ratio 13 : 8 : 1, their areas are inthe ratio 169 : 64 : 1. The area of triangle ABC is 1/2 (12)(5) = 30, so the areas of triangle NBK and triangle AMJ are (64/169) (30) and (1/169)(30), respectively. Thus the area of pentagon CMJKN is (1 -  64/169 - 1/169)(30)= $\boxed{\mathbf{(D)}240/13}$

 Credit to \href{https://artofproblemsolving.comhttp://billingswest.billings.k12.mt.us/math/AMC\%201012/AMC\%2012\%20work\%20sheets/2004\%20AMC\%2012B\%20ws-15.pdf}{http://billingswest.billings.k12.mt.us/math/AMC\%201012/AMC\%2012\%20work\%20sheets/2004\%20AMC\%2012B\%20ws-15.pdf} for Solution 3.

 



\textbf{Solution 5}

We fake-solve. Draw a perpendicular line to side $CB$ that goes through $N$. Similarly, draw a perpendicular line to side $AC$ that goes through $M$. Let them intersect at $O$. It follows that quadrilateral $MONC$ is a square with side $4$ and thus area $16$. Now, draw a perpendicular from $J$ to $MO$ and let them intersect at $Z$. There is a obvious pair of congruent triangles. Fill in the gap. Hence, we see that the area of the pentagon is definitely greater than $16$. How much greater? Well, if we let the intersection of $AB$ and $NO$ be $X$, we can approximate that $\triangle NXK$ has an area greater than $2$. So we place bets on $\boxed{\mathbf{(D)}240/13}$ which is correct!

 This would be much easier to show and more convincing with a diagram but idk how to upload GeoGebra to AoPSwiki.

 Solution by franzliszt

 


	\item (\textit{AMC 12 12B 2004}) The two digits in Jack's age are the same as the digits in Bill's age, but in reverse order. In five years Jack will be twice as old as Bill will be then. What is the difference in their current ages?


 $\mathrm{(A) \ } 9 \qquad \mathrm{(B) \ } 18 \qquad \mathrm{(C) \ } 27 \qquad \mathrm{(D) \ } 36\qquad \mathrm{(E) \ } 45$

 \textbf{Answer: }B



\textbf{Solution 1}

If Jack's current age is $\overline{ab}=10a+b$, then Bill's current age is $\overline{ba}=10b+a$.

 In five years, Jack's age will be $10a+b+5$ and Bill's age will be $10b+a+5$. 

 We are given that $10a+b+5=2(10b+a+5)$. 

 
Thus $8a=19b+5 \Rightarrow a=\dfrac{19b+5}{8}$. 

 For $b=1$ we get $a=3$. For $b=2$ and $b=3$ the value $\frac{19b+5}8$ is not an integer, and for $b\geq 4$, $a$ is more than $9$. Thus the only solution is $(a,b)=(3,1)$, and the difference in ages is $31-13=\boxed{\mathrm{(B)\ }18}$.

 



\textbf{Solution 2}

Age difference does not change in time. Thus in five years Bill's age will be equal to their age difference.

 The age difference is $(10a+b)-(10b+a)=9(a-b)$, hence it is a multiple of $9$. Thus Bill's current age modulo $9$ must be $4$.

 Thus Bill's age is in the set $\{13,22,31,40,49,58,67,76,85,94\}$.

 As Jack is older, we only need to consider the cases where the tens digit of Bill's age is smaller than the ones digit. This leaves us with the options $\{13,49,58,67\}$. 

 Checking each of them, we see that only $13$ works, and gives the solution $31-13=\boxed{\mathrm{(B)\ }18}$.

 


	\item (\textit{AMC 12 12B 2004}) A \href{https://artofproblemsolving.com/wiki/index.php?title=Function}{function} $f$ is defined by $f(z) = i\overline{z}$, where $i=\sqrt{-1}$ and $\overline{z}$ is the \href{https://artofproblemsolving.com/wiki/index.php?title=Complex\_conjugate}{complex conjugate} of $z$. How many values of $z$ satisfy both $|z| = 5$ and $f(z) = z$?


 $\mathrm{(A)}\ 0 \qquad\mathrm{(B)}\ 1 \qquad\mathrm{(C)}\ 2  \qquad\mathrm{(D)}\ 4 \qquad\mathrm{(E)}\ 8$

 \textbf{Answer: }C



\textbf{Solution 1}

Let $z = a+bi$, so $\overline{z} = a-bi$. By definition, $z = a+bi = f(z) = i(a-bi) = b+ai$, which implies that all solutions to $f(z) = z$ lie on the line $y=x$ on the complex plane. The graph of $|z| = 5$ is a \href{https://artofproblemsolving.com/wiki/index.php?title=Circle}{circle} centered at the origin, and there are $2 \Rightarrow \mathrm{(C)}$ intersections.

 



\textbf{Solution 2}

Let $z=a+bi$, like above. Therefore, $z = a+bi = i\overline{z} = i(a-bi) = ai+b$. We move some terms around to get $bi-b = ai-a$. We factor: $b(i-1) = a(i-1)$. We divide out the common factor to see that $b = a$. Next we put this into the definition of $|z| = a^2 + b^2 = a^2 + a^2 = 2a^2 = 25$. Finally, $a = \pm\sqrt{\frac{25}{2}}$, and $a$ has two solutions.

 



\textbf{Solution 3}

Via trig, let $z=r(\cos(x)+i\sin(x))$. Since we wish for $|z|=5$, let $r=5$. We want to find $x\in [0, 2\pi)$ such that $f(z)=i\overline{z}$. Note that $\overline{z}=5(\cos(x)-i\sin(x))$, meaning $i\overline{z}=5(i\cos(x)+\sin(x))$. Thus, we have: $5(\cos(x)+i\sin(x))=5(i\cos(x)+\sin(x))$. Grouping the real and imaginary values, we get that $\cos(x)=\sin(x)$. Note that this has two solutions in $[0, 2\pi)$. Thus, there are $2$ valid solutions.

 


	\item (\textit{AMC 12 12B 2004}) For some real numbers $a$ and $b$, the \href{https://artofproblemsolving.com/wiki/index.php?title=Equation}{equation}

 \[8x^3 + 4ax^2 + 2bx + a = 0\]
has three distinct positive roots. If the sum of the base-$2$ \href{https://artofproblemsolving.com/wiki/index.php?title=Logarithm}{logarithms} of the roots is $5$, what is the value of $a$? 


 $\mathrm{(A)}\ -256 \qquad\mathrm{(B)}\ -64  \qquad\mathrm{(C)}\ -8 \qquad\mathrm{(D)}\ 64  \qquad\mathrm{(E)}\ 256$

 \textbf{Answer: }A



\textbf{Solution}

Let the three roots be $x_1,x_2,x_3$. 
 \[\log_2 x_1 + \log_2 x_2 + \log_2 x_3 = \log_2 x_1x_2x_3= 5 \Longrightarrow x_1x_2x_3 = 32\]By \href{https://artofproblemsolving.com/wiki/index.php?title=Vieta\%27s\_Formulas}{Vieta’s formulas},
 \[8(x-x_1)(x-x_2)(x-x_3) = 8x^3 + 4ax^2 + 2bx + a\]gives us that $a = -8x_1x_2x_3 = -256 \Rightarrow \mathrm{(A)}$.

 


	\item (\textit{AMC 12 12B 2004}) Points $A$ and $B$ are on the parabola $y=4x^2+7x-1$, and the origin is the midpoint of $AB$. What is the length of $AB$?


 $\mathrm{(A)}\ 2\sqrt5 \qquad \mathrm{(B)}\ 5+\frac{\sqrt2}{2} \qquad \mathrm{(C)}\ 5+\sqrt2 \qquad \mathrm{(D)}\ 7 \qquad \mathrm{(E)}\ 5\sqrt2$

 \textbf{Answer: }E



\textbf{Solution}

Let the coordinates of $A$ be $(x_A,y_A)$. As $A$ lies on the parabola, we have $y_A=4x_A^2+7x_A-1$.As the origin is the midpoint of $AB$, the coordinates of $B$ are $(-x_A,-y_A)$.We need to choose $x_A$ so that $B$ will lie on the parabola as well. In other words, we need $-y_A = 4(-x_A)^2 + 7(-x_A) - 1$.

 Substituting for $y_A$, we get: $-4x_A^2 - 7x_A + 1 = 4(-x_A)^2 + 7(-x_A) - 1$.

 This simplifies to $8x_A^2 - 2 = 0$, which solves to $x_A = \pm 1/2$. Both roots lead to the same pair of points: $(1/2,7/2)$ and $(-1/2,-7/2)$. Their distance is $\sqrt{ 1^2 + 7^2 } = \sqrt{50} = \boxed{5\sqrt2}$.

 \textbf{Alternate Solution} Let the coordinates of $A$ and $B$ be $(x_A, y_A)$ and $(x_B, y_B)$, respectively. Since the median of the points lies on the origin, $x_A + x_B = y_A + y_B = 0$ and expanding $y_A + y_B$, we find:
 \[4x_A^2 + 7x_A - 1 + 4x_B^2 + 7x_B - 1 = 0\]
 \[4(x_A^2 + x_B^2) + 7(x_A + x_B) = 2\]
 \[x_A^2 + x_B^2 = \frac{1}{2}.\]

 It also follows that $(x_A + x_B)^2 = 0$. Expanding this, we find:
 \[x_A^2 + 2x_A x_B + x_B^2 = 0\]
 \[\frac{1}{2} + 2x_A x_B = 0\]
 \[x_A x_B = -\frac{1}{4}.\]

 To find the distance between the points, $\sqrt{(x_A - x_B)^2 + (y_A - y_B)^2}$ must be found. Expanding $y_A - y_B$:
 \[y_A - y_B = 4x_A^2 + 7x_A - 1 - 4x_B^2 - 7x_B + 1\]
 \[= 4(x_A^2 - x_B^2) + 7(x_A - x_B)\]
 \[= 4(x_A + x_B)(x_A - x_B) + 7(x_A - x_B)\]
 \[= 7(x_A - x_B)\]we find the distance to be $\sqrt{50(x_A - x_B)^2}$. Expanding this yields $5\sqrt{2(x_A^2 + x_B^2 - 2x_A x_B)} = \boxed{5\sqrt{2}}$.

 


	\item (\textit{AMC 12 12B 2004}) A truncated \href{https://artofproblemsolving.com/wiki/index.php?title=Cone}{cone} has horizontal bases with radii $18$ and $2$. A \href{https://artofproblemsolving.com/wiki/index.php?title=Sphere}{sphere} is tangent to the top, bottom, and lateral surface of the truncated cone. What is the \href{https://artofproblemsolving.com/wiki/index.php?title=Radius}{radius} of the sphere? 


 $\mathrm{(A)}\ 6 \qquad\mathrm{(B)}\ 4\sqrt{5} \qquad\mathrm{(C)}\ 9 \qquad\mathrm{(D)}\ 10 \qquad\mathrm{(E)}\ 6\sqrt{3}$

 \textbf{Answer: }A



\textbf{Solution}

Consider a \href{https://artofproblemsolving.com/wiki/index.php?title=Trapezoid}{trapezoid} (label it $ABCD$ as follows) cross-section of the truncate cone along a diameter of the bases:

 
                    \begin{figure}[H]
                        \centering
                        \includegraphics[width=0.5\textwidth]{Solution/AMC_12_12B_2004_19_Solution_Figure_1.png}
                    \end{figure}
                    

 Above, $E,F,$ and $G$ are points of \href{https://artofproblemsolving.com/wiki/index.php?title=Tangent\_(geometry)}{tangency}. By the Two Tangent Theorem, $BF = BE = 18$ and $CF = CG = 2$, so $BC = 20$. We draw $H$ such that it is the foot of the altitude $\overline{HD}$ to $\overline{AB}$: 

 
                    \begin{figure}[H]
                        \centering
                        \includegraphics[width=0.5\textwidth]{Solution/AMC_12_12B_2004_19_Solution_Figure_2.png}
                    \end{figure}
                    

 By the \href{https://artofproblemsolving.com/wiki/index.php?title=Pythagorean\_Theorem}{Pythagorean Theorem}, 
 \[r = \frac{DH}2 = \frac{\sqrt{20^2 - 16^2}}2 = 12\]Therefore, the answer is $\boxed{{A (6)}}.$



	\item (\textit{AMC 12 12B 2004}) Each face of a \href{https://artofproblemsolving.com/wiki/index.php?title=Cube}{cube} is painted either red or blue, each with probability $1/2$. The color of each face is determined independently. What is the probability that the painted cube can be placed on a horizontal surface so that the four vertical faces are all the same color?


 $\textbf{(A)}\ \frac14 \qquad \textbf{(B)}\ \frac {5}{16} \qquad \textbf{(C)}\ \frac38 \qquad \textbf{(D)}\ \frac {7}{16} \qquad \textbf{(E)}\ \frac12$

 \textbf{Answer: }B



\textbf{Solution}

There are $2^6$ possible colorings of the cube. Consider the color that appears with greater frequency. The property obviously holds true if $5$ or $6$ of the faces are colored the same, which for each color can happen in $6 + 1 = 7$ ways. If $4$ of the faces are colored the same, there are $3$ possible cubes (corresponding to the $3$ possible ways to pick pairs of opposite faces for the other color). If $3$ of the faces are colored the same, the property obviously cannot be satisfied. Thus, there are a total of $2(7 + 3) = 20$ ways for this to occur, and the desired probability is $\frac{20}{2^6} = \frac{5}{16}\ \mathbf{(B)}$.

 



\textbf{Solution 2}

As with the previous solution, we will simply count the number of valid faces and divide by $2^6$. We say that a cube has a ring of color when four vertical faces are all the same color. If we fix the cube in space, there can be $3$ different rings (in each of the $3$ axis), each of which can take on 2 different colors. This leaves us to choose the remaining two faces (which are opposite one another) in $2^2=4$ ways. However, we have counted coloring that has each face the same color $3$ times for both red and blue. This means that we must subtract the 4 overcounts to get $2*3*2^2-4=20$. Dividing by $2^6$ total colorings, we get $20/2^6=\frac{5}{16}\mathbf{(B)}$.~bambithenambi

 


	\item (\textit{AMC 12 12B 2004}) The graph of $2x^2 + xy + 3y^2 - 11x - 20y + 40 = 0$ is an ellipse in the first quadrant of the $xy$-plane. Let $a$ and $b$ be the maximum and minimum values of $\frac yx$ over all points $(x,y)$ on the ellipse. What is the value of $a+b$?


 $\mathrm{(A)}\ 3 \qquad\mathrm{(B)}\ \sqrt{10} \qquad\mathrm{(C)}\ \frac 72 \qquad\mathrm{(D)}\ \frac 92 \qquad\mathrm{(E)}\ 2\sqrt{14}$

 \textbf{Answer: }C



\textbf{Solution}


                    \begin{figure}[H]
                        \centering
                        \includegraphics[width=0.5\textwidth]{Solution/AMC_12_12B_2004_21_Solution_Figure_1.png}
                    \end{figure}
                    
$\frac yx$ represents the slope of a line passing through the origin. It follows that since a line $y = mx$ intersects the ellipse at either $0, 1,$ or $2$ points, the minimum and maximum are given when the line $y = mx$ is a tangent, with only one point of intersection. Substituting,
 \[2x^2 + x(mx) + 3(mx)^2 - 11x - 20(mx) + 40 = 0\]Rearranging by the degree of $x$,
 \[(3m^2 + m + 2)x^2 - (20m + 11)x + 40 = 0\]Since the line $y=mx$ is tangent to the ellipse, we want the discriminant, 
 \[(20m+11)^2 - 4\cdot 40 \cdot (3m^2 + m + 2) = -80m^2 + 280m - 199\]to be equal to $0$. We want $a+b$, which is the sum of the roots of the above quadratic. By \href{https://artofproblemsolving.com/wiki/index.php?title=Vieta\%E2\%80\%99s\_formulas&action=edit&redlink=1}{Vieta’s formulas}, that is $\frac{280}{80} = \frac{7}{2} \Rightarrow \mathrm{(C)}$.

 


	\item (\textit{AMC 12 12B 2004}) The square


 $\begin{tabular}{|c|c|c|} \hline 50 & \textit{b} & \textit{c} \\ \hline \textit{d} & \textit{e} & \textit{f} \\ \hline \textit{g} & \textit{h} & 2 \\ \hline \end{tabular}$

 is a multiplicative magic square. That is, the product of the numbers in each row, column, and diagonal is the same. If all the entries are positive integers, what is the sum of the possible values of $g$?


 $\textbf{(A)}\ 10 \qquad \textbf{(B)}\ 25 \qquad \textbf{(C)}\ 35 \qquad \textbf{(D)}\ 62 \qquad \textbf{(E)}\ 136$

 \textbf{Answer: }C



\textbf{Solution 1}

All the unknown entries can be expressed in terms of $b$.Since $100e = beh = ceg = def$, it follows that $h = \frac{100}{b}, g = \frac{100}{c}$,and $f = \frac{100}{d}$. Comparing rows $1$ and $3$ then gives$50bc = 2 \cdot \frac{100}{b} \cdot \frac{100}{c}$,from which $c = \frac{20}{b}$.Comparing columns $1$ and $3$ gives$50d \cdot \frac{100}{c}= 2c \cdot \frac{100}{d}$,from which $d = \frac{c}{5} = \frac{4}{b}$.Finally, $f = 25b, g = 5b$, and $e = 10$. All the entries are positive integersif and only if $b = 1, 2,$ or $4$. The corresponding values for $g$ are $5, 10,$ and$20$, and their sum is $\boxed{\mathbf{(C)}35}$.

 



\textbf{Solution 2}

We know because this is a multiplicative magic square that each of the following are equal to each other:$100e=ceg=50dg=beh=2cf=50bc=def=2gh$

 From this we know that $50dg=2hg$, thus $h=25d$.Thus $beh=be(25d)$ and $be(25d)=100e$. Thus $b=\frac{4}{d}$From this we know that $50bc=(50)(\frac{4}{d})(c)=50dg$. Thus $c=\frac{d^2g}{4}$.Now we know from the very beginning that $100e=ceg$ or $100=cg$ or $100=\frac{d^2g}{4}(g)$ or $\frac{d^2g^2}{4}$. Rearranging the equation $100=\frac{d^2g^2}{4}$ we have $(d^2)(g^2)=400$ or $dg=20$ due to $d$ and $g$ both being positive. Now that $dg=20$ we find all pairs of positive integers that multiply to $20$. There is $(d,g)= (20,1);(10,2);(5,4);(4,5);(2,10);(1,20)$. Now we know that $b=\frac{4}{d}$ and b has to be a positive integer. Thus $d$ can only be $1$, $2$, or $4$. Thus $g$ can only be $20$,$10$,or $5$. Thus sum of $20+10+5$ = $35$. The answer is $\boxed{\mathbf{(C)}35}$.

 


	\item (\textit{AMC 12 12B 2004}) The \href{https://artofproblemsolving.com/wiki/index.php?title=Polynomial}{polynomial} $x^3 - 2004 x^2 + mx + n$ has \href{https://artofproblemsolving.com/wiki/index.php?title=Integer}{integer} coefficients and three distinct positive zeros. Exactly one of these is an integer, and it is the sum of the other two.  How many values of $n$ are possible?


 $\mathrm{(A)}\ 250,000 \qquad\mathrm{(B)}\ 250,250 \qquad\mathrm{(C)}\ 250,500 \qquad\mathrm{(D)}\ 250,750 \qquad\mathrm{(E)}\ 251,000$

 \textbf{Answer: }C



\textbf{Solution 1}

Let the roots be $r,s,r + s$, and let $t = rs$. Then

 $(x - r)(x - s)(x - (r + s))$ $= x^3 - (r + s + r + s) x^2 + (rs + r(r + s) + s(r + s))x - rs(r + s) = 0$

 and by matching coefficients, $2(r + s) = 2004 \Longrightarrow r + s = 1002$. Then our polynomial looks like 
 \[x^3 - 2004x^2 + (t + 1002^2)x - 1002t = 0\]and we need the number of possible products $t = rs = r(1002 - r)$. Because $m=t+1002^2$ is an integer, we also note that $t$ must be an integer. 

 Since $r > 0$ and $t > 0$, it follows that $0 < t = r(1002-r) < 501^2 = 251001$, with the endpoints not achievable because the roots must be distinct and positive. Because neither $r$ nor $1002-r$ can be an integer, there are $251,000 - 500 = \boxed{\textbf{(C) } 250,500}$ possible values of $n = -1002t$.

 


 



\textbf{Solution 2}

Letting the roots be $r$, $s$, and $t$, where $t = r+s$, we see that by Vieta's Formula's, $2004 = r+s+t = t + t = 2t$, and so $t = 1002$. Therefore, $x-1002$ is a factor of $x^3 - 2004x^2 + mx + n$. Letting $x = 0$ gives that $1002 \mid n$ because $x - 1002 \mid x^3 - 2004x^2 + mx + n$. Letting $n = -1002a$ and noting that $x^3 - 2004x^2 + mx + n = (x-1002)(x^2 - bx + a)$ for some $b$, we see that $b$ is the sum of the roots of $x^2 - bx + a$, $r$ and $s$, and so $b = 1002$. Now, we have that $x^2 - 1002x + a$ has roots $r$ and $s$, and we wish to find the number of possible values of $a$. By the quadratic formula, we see that 
 \[\frac{1002 \pm \sqrt{1002^2 - 4a}}{2} = 501 \pm \sqrt{501^2 - a}\] are the two values of noninteger positive real numbers $r$ and $s$, neither of which is equal to $1002$. This information gives us that $0 < 501^2 - a < 501^2$, and so since $501^2 - a$ is evidently not a square, we have $501^2 - 1 - 500 = 251,001 - 500 - 1 = \boxed{\textbf{(C) } 250,500}$ possible values of $n = 1002a$.

 


	\item (\textit{AMC 12 12B 2004}) In $\triangle ABC$, $AB = BC$, and $\overline{BD}$ is an \href{https://artofproblemsolving.com/wiki/index.php?title=Altitude}{altitude}. Point $E$ is on the extension of $\overline{AC}$ such that $BE = 10$. The values of $\tan \angle CBE$, $\tan \angle DBE$, and $\tan \angle ABE$ form a \href{https://artofproblemsolving.com/wiki/index.php?title=Geometric\_progression}{geometric progression}, and the values of $\cot \angle DBE,$ $\cot \angle CBE,$ $\cot \angle DBC$ form an \href{https://artofproblemsolving.com/wiki/index.php?title=Arithmetic\_progression}{arithmetic progression}. What is the area of $\triangle ABC$?


 
                    \begin{figure}[H]
                        \centering
                        \includegraphics[width=0.5\textwidth]{Problem/AMC_12_12B_2004_24_Problem_Figure_1.png}
                    \end{figure}
                    
$\mathrm{(A)}\ 16 \qquad\mathrm{(B)}\ \frac {50}3 \qquad\mathrm{(C)}\ 10\sqrt{3} \qquad\mathrm{(D)}\ 8\sqrt{5} \qquad\mathrm{(E)}\ 18$

 \textbf{Answer: }B



\textbf{Solution}

Let $\alpha = DBC$. Then the first condition tells us that
 \[\tan^2 DBE = \tan(DBE - \alpha)\tan(DBE + \alpha) = \frac {\tan^2 DBE - \tan^2 \alpha}{1 - \tan ^2 DBE \tan^2 \alpha},\]and multiplying out gives us $(\tan^4 DBE - 1) \tan^2 \alpha = 0$. Since $\tan\alpha \neq 0$, we have $\tan^4 DBE = 1 \Longrightarrow \angle DBE = 45^{\circ}$. 

 The second condition tells us that $2\cot (45 - \alpha) = 1 + \cot \alpha$. Expanding, we have $1 + \cot \alpha = 2\left[\frac {\cot \alpha + 1}{\cot \alpha - 1}\right] \Longrightarrow (\cot \alpha - 3)(\cot \alpha + 1) = 0$. Evidently $\cot \alpha \neq - 1$, so we get $\cot \alpha = 3$.

 Now $BD = 5\sqrt {2}$ and $AC = \frac {2BD} {\cot \alpha} = \frac {10\sqrt {2}}{3}$. Thus, $[ABC] = \frac {1}{2} \cdot 5\sqrt {2} \cdot \frac {10\sqrt {2}}{3} = \frac {50}{3}\ \mathrm{(B)}$.

 


	\item (\textit{AMC 12 12B 2004}) Given that $2^{2004}$ is a $604$-\href{https://artofproblemsolving.com/wiki/index.php?title=Digit}{digit} number whose first digit is $1$, how many \href{https://artofproblemsolving.com/wiki/index.php?title=Element}{elements} of the \href{https://artofproblemsolving.com/wiki/index.php?title=Set}{set} $S = \{2^0,2^1,2^2,\ldots ,2^{2003}\}$ have a first digit of $4$? 


 $\mathrm{(A)}\ 194 \qquad\mathrm{(B)}\ 195 \qquad\mathrm{(C)}\ 196 \qquad\mathrm{(D)}\ 197 \qquad\mathrm{(E)}\ 198$

 \textbf{Answer: }B



\textbf{Solution 1}

Given $n$ digits, there must be exactly one power of $2$ with $n$ digits such that the first digit is $1$. Thus $S$ contains $603$ elements with a first digit of $1$. For each number in the form of $2^k$ such that its first digit is $1$, then $2^{k+1}$ must either have a first digit of $2$ or $3$, and $2^{k+2}$ must have a first digit of $4,5,6,7$. Thus there are also $603$ numbers with first digit $\{2,3\}$ and $603$ numbers with first digit $\{4,5,6,7\}$. By using \href{https://artofproblemsolving.com/wiki/index.php?title=Complementary\_counting}{complementary counting}, there are $2004 - 3 \times 603 = 195$ elements of $S$ with a first digit of $\{8,9\}$. Now, $2^k$ has a first digit of $\{8,9\}$ \href{https://artofproblemsolving.com/wiki/index.php?title=Iff}{if and only if} the first digit of $2^{k-1}$ is $4$, so there are $\boxed{195} \Rightarrow \mathrm{(B)}$ elements of $S$ with a first digit of $4$.

 



\textbf{Solution 2}

We can make the following chart for the possible loops of leading digits:
 \[1 \rightarrow 2 \rightarrow 4 \rightarrow 8 \rightarrow 1\]
 \[1 \rightarrow 2 \rightarrow 4 \rightarrow 9 \rightarrow 1\]
 \[1 \rightarrow 2 \rightarrow 5 \rightarrow 1\]
 \[1 \rightarrow 3 \rightarrow 6 \rightarrow 1\]
 \[1 \rightarrow 3 \rightarrow 7 \rightarrow 1\]

 
Thus each loop from $1 \rightarrow 1$ can either have $3$ or $4$ numbers. Let there be $x$ of the sequences of $3$ numbers, and let there be $y$ of the sequences of $4$ numbers. We note that a $4$ appears only in the loops of $4$, and also we are given that $2^{2004}$ has $604$ digits.
 \[3x+4y=2004\]
 \[x+y=603\]Solving gives $x = 408$ and $y = 195$, thus the answer is $(B)$.

 


\end{enumerate}